The energy certificate developed in Section~\ref{sec:4.A} rests on a prediction of remaining voltage, whose error the online calibration keeps within an interval.
We assume the interval exists and defer its construction to Section~\ref{sec:4.A}.

\begin{assumption}[Calibrated proxy] \label{as:proxy}
The proxy's per-step prediction error is calibrated online: for the residual $s_t$ of Eq.~\eqref{eq:residual} between predicted and realized voltage drop over an interval $\Delta t$, adaptive conformal inference returns a half-width $q_t$ whose long-run exceedance frequency $\lim_{T\to\infty} \frac{1}{T} \sum_{t=1}^{T} \mathbb{1}[s_t > q_t]$ equals $\delta$. 
The proxy error is homogeneous across the admissible input set, so a half-width calibrated on realized inputs bounds the error of the nominal-policy prediction used in $\hat{c}_\theta$ and $\hat c_\theta^{g_i}$.
\end{assumption}

\begin{assumption}[External goal allocation] \label{as:allocator}
Each robot is assigned a goal $p_i^g$, either a task location or a station in $\mathcal{C}$, by an external allocator treated as a black box. 
Goals are pairwise non-conflicting, $\|p_i^g - p_j^g\| > 2r$, and every assignment is individually survivable: under the calibrated proxy, robot $i$ can reach $p_i^g$ and thereafter a station with voltage no less than $V_{\min} + \varepsilon_i$; otherwise it is redirected to a station. 
The allocator certifies only this per-robot feasibility. 
It does not compare robots, and the relative ranking that governs yielding is formed at runtime in Section~\ref{sec:4.B}.
\end{assumption}

\begin{assumption}[Feasible initial conditions]\label{as:init}
At $t=0$, the robots are mutually safe, $\|p_i(0) - p_j(0)\| > 2r$ and $\mathrm{dist}(p_i(0), \mathcal{O}) > r$, and each can reach some station with margin $V_{\min} + \varepsilon_i$ under the calibrated proxy.
\end{assumption}

\begin{problem}[Cooperative Safe Control with Energy-Guided Precedence]\label{prob:main}
We seek a distributed policy $u_i=\pi_i\big(z_i, \{z_j\}_{j\in\mathcal{N}_i}, p_i^g\big)$ under which the closed-loop trajectories of Eq.~\eqref{eq:augmented_dynamics}, with resets, satisfy the following for every robot $i$ and all $t \geq 0$:
\begin{enumerate}
    \item \textit{Collision safety:} $\|p_i(t) - p_j(t)\| \geq 2r$ for all $j\neq i$, and $\mathrm{dist}(p_i(t), \mathcal{O})\ge r$.
    \item \textit{Energy sufficiency:} robot $i$ retains enough voltage to reach a station at all times, certified through the reachability set of Section~\ref{sec:4.A} rather than through $\{V_i \geq V_{\min}\}$, despite the uncertainty in $\theta$.
    \item \textit{Liveness:} $\inf_{t\geq 0} \|p_i(t) - p_i^g
    \|=0$, i.e., each robot eventually reaches its assigned goal.
\end{enumerate}
\end{problem}

Collision safety (i) is enforced unconditionally.
Energy sufficiency (ii) holds throughout the mission and the subsequent return, so a robot that has completed its task is still certified to reach a station. 
Among policies meeting (i)–(iii), we prefer those that yield according to energy.
When two robots must give way, the one with less voltage to spare deviates less, and its neighbor absorbs the effort. 
This yielding is a preference layered on top of the safety guarantees, upheld whenever it does not conflict with (i).
We assume a nominal controller $u_i^\mathrm{nom}$ that achieves (iii) but not necessarily (i) or (ii). 
The task is a distributed policy that enforces the latter two while deviating minimally from $u_i^\mathrm{nom}$, with the deviation steered by precedence toward the robots that can afford it.